
\section{The bridge to the main paper: block decomposition}
\label{sec:blocks}

This section connects the theory above to the indicator of the main paper. It
establishes the identity behind the ``contagion map'' and shows where the
dependence channel sits within the structural index.

\subsection{The identity}

Partition the $N$ series into $B$ blocks (for instance the CISS segments: money
market, bonds, equities, financial intermediaries, foreign exchange), and write
$R_{bb}$ for the diagonal blocks of the correlation matrix $R$ and
$R_D=\operatorname{diag}(R_{11},\dots,R_{BB})$.

\begin{proposition}[Within/between-block decomposition]
\label{prop:blocks}
\begin{equation}
\label{eq:blocks}
\underbrace{-\tfrac12\log\det R}_{\text{total dependence}}
\;=\;\sum_{b=1}^{B}\underbrace{\left(-\tfrac12\log\det R_{bb}\right)}_{\text{within block }b}
\;+\;\underbrace{\left(-\tfrac12\log\det\big(R_D^{-1/2}RR_D^{-1/2}\big)\right)}_{\text{between blocks}} ,
\end{equation}
and each of the $B+1$ terms is non-negative. The between-block term is the
Gaussian multi-information among blocks,
$\KL{\mathcal N(0,R)}{\bigotimes_b\mathcal N(0,R_{bb})}$; it is zero if and only
if the blocks are pairwise uncorrelated.
\end{proposition}

\begin{proof}
$R=R_D^{1/2}\,\big(R_D^{-1/2}RR_D^{-1/2}\big)\,R_D^{1/2}$ and the determinant is
multiplicative, so $\log\det R=\log\det R_D+\log\det(R_D^{-1/2}RR_D^{-1/2})$,
with $\log\det R_D=\sum_b\log\det R_{bb}$. Each $R_{bb}$ is a correlation matrix,
so $-\tfrac12\log\det R_{bb}\ge0$ by Hadamard's inequality
(Appendix~\ref{app:gauss}). The matrix $\tilde R=R_D^{-1/2}RR_D^{-1/2}$ has
identity diagonal blocks, hence a unit diagonal, and Hadamard applies again.
Finally Proposition~\ref{prop:klgauss} with $\mu_1=\mu_0$, $\Sigma_1=R$,
$\Sigma_0=R_D$ gives
$\tfrac12[\operatorname{tr}(R_D^{-1}R)-N-\log\det(R_D^{-1}R)]$, and
$\operatorname{tr}(R_D^{-1}R)=\sum_b\operatorname{tr}(I)=N$.
\end{proof}

\paragraph{Two blocks: canonical correlations.} For $B=2$, writing
$C=R_{11}^{-1/2}R_{12}R_{22}^{-1/2}$ and $\rho_1,\dots,\rho_k$ for its singular
values (the canonical correlations between the two blocks),
$\det\tilde R=\det(I-CC^\top)=\prod_j(1-\rho_j^2)$ and the between-block term
equals $-\tfrac12\sum_j\log(1-\rho_j^2)$. This is the Gaussian mutual
information between the two segments, the sum of the information carried by each
canonical pair. For $B>2$, identity~\eqref{eq:blocks} can be applied in nested
fashion and lets one attribute contagion to pairs or groups of segments by the
chain rule.

\subsection{Where the dependence channel sits in the structural index}

The chain rule of Appendix~\ref{app:index} writes $\JJ=\sum_i\Dun(x_i)+I(x)$.
Replacing the multi-information $I(x)$ by its Gaussian part $-H^{dep}$ and
applying~\eqref{eq:blocks}:
\begin{equation}
\label{eq:jblocks}
\JJ\;\approx\;\underbrace{\sum_i\Dun(x_i)}_{\text{shape}}
\;+\;\underbrace{\sum_b\left(-\tfrac12\log\det R_{bb}\right)}_{\text{within}}
\;+\;\underbrace{\left(-\tfrac12\log\det\tilde R\right)}_{\text{between}} .
\end{equation}
The main paper keeps the last two terms and drops the first, which is the one
the null reproduces (Section~\ref{sec:program}). The entropic CISS is thus the
structural index stripped of its shape channel, and its decomposition is exact,
with no approximation, because it concerns only the Gaussian part.

\subsection{Three properties useful for the indicator}

\paragraph{Scale invariance.} The terms of~\eqref{eq:blocks} depend only on
$R$: they are invariant to any series-by-series rescaling, which allows markets
with very different volatilities to be aggregated without ad hoc weighting. The
\citet{forbes2002} devolatilisation therefore acts not on the formula but on the
bias of the \emph{estimator} of $R$ under heteroskedasticity.

\paragraph{Rank version.} Replacing $R$ by the correlation of normal scores
$\Phi^{-1}(\hat F_i(x_i))$ estimates the multi-information of a Gaussian copula.
The quantity becomes invariant to any monotone transformation of the marginals,
which explains its robustness to tails observed in Phase~A.

\paragraph{Normalisation in $N$.} For a constant correlation $\rho$,
$\det R=(1-\rho)^{N-1}\big(1+(N-1)\rho\big)$, so $-H^{dep}$ grows roughly linearly
in $N$. Dividing by $N-1$, or reporting the equivalent constant correlation
$\rho^\star$ that gives the same determinant, makes the indicator comparable
across panels. For small $\rho^\star$,
$-H^{dep}\approx\tfrac14N(N-1)\rho^{\star2}$.

\section{Monitoring indicators and insurance solvency}
\label{sec:monitoring}

\paragraph{Two indicators from v1.} The jump/diffusion split delivered two daily
quantities: the standardised positive-jump component of $\JJ$ (how much
information arrived today) and $\JJ_t$ relative to its rolling maximum (how much
has not yet been absorbed). Neither improves the volatility forecast (H7), and
both rest on statistics the null does not identify. They are not carried over.

\paragraph{Insurance.} Three links remain valid in principle. An entropy-ball
scenario produces a complete stressed law, a natural object for the ORSA. The
decomposition of classical bundles bears on the Solvency~II standard formula,
whose diversification credit rests on a fixed correlation matrix. And entropic
risk measures are coherent \citep{ahmadijavid2012,artzner1999}. Nothing has been
tested on liabilities, however, and the short, low-frequency panels typical of
insurance data would strain every estimator used here.

\section{Open questions}
\label{sec:open}

\begin{enumerate}
\item \textbf{Self-exciting arrivals.} The jump clustering seen in the v1 event
study contradicts the Poisson measure in~\eqref{eq:ou}. A Hawkes process is the
next specification.
\item \textbf{Positivity.} $\JJ\ge0$ whereas an Ornstein--Uhlenbeck process with
a Brownian term can turn negative. One should either drop $\sigma_\JJ$ (the
pure form of \citealt{bns2001}) or model $\log\JJ$.
\item \textbf{Non-linear tail dependence.} $\JJ$ and the dependence channel see
only the Gaussian part of the multi-information. Estimating the full copula
multi-information (for example by nearest neighbours on ranks) would measure
what is missing, at the price of a much higher estimation variance.
\item \textbf{Tail of the information flow.} Beyond about twenty years, the
quantiles of $I^{(h)}$ are extrapolated. A generalised Pareto fit would be the
disciplined correction.
\item \textbf{Minimal lift and published scenarios.}
Proposition~\ref{prop:lift} provides a yardstick: for any published scenario,
the ratio of its ambient price to the minimal price $g(u,v)$ of its projection
measures how much information it spends outside the direction that matters to
the portfolio. This extends the audit of the main paper directly.
\end{enumerate}
